\documentclass[12pt]{article}
\usepackage[utf8]{inputenc}
\usepackage[T1]{fontenc}
\usepackage[margin=1in]{geometry}
\usepackage{setspace}
\usepackage{xcolor}
\usepackage[colorlinks=true, urlcolor=blue!60!black, linkcolor=blue!60!black]{hyperref}
\setlength{\parindent}{2em}
\setlength{\parskip}{1.25em}
\renewcommand{\baselinestretch}{1.0}
\begin{document}
\begin{center}
{\LARGE \textit{Cohen v.\ California}, 403 U.S.\ 15 (1971)\par}
\vspace{0.8em}
{\large Case Brief --- Four-Page Edition\par}
\end{center}
\vspace{1em}

\section{Memory Jogger}

A 19-year-old wore a jacket reading ``Fuck the Draft'' in a Los Angeles courthouse corridor. The Court held 5--4 that California could not punish the jacket's political message under its general breach-of-the-peace statute, because it was neither obscenity nor fighting words and unwilling viewers could avert their eyes (\href{https://supreme.justia.com/cases/federal/us/403/15/}{Justia}).

\section{Detailed Case Facts}

On April 26, 1968, nineteen-year-old Paul Robert Cohen walked through a corridor of the Los Angeles County Courthouse outside a division of the Municipal Court wearing a jacket bearing the plainly visible words ``Fuck the Draft.'' Women and children were present. When Cohen entered a courtroom, he removed the jacket and held it folded over his arm; a police officer asked the judge to hold him in contempt, the judge declined, and the officer arrested Cohen when he stepped back into the corridor. Cohen testified that he wore the jacket to show the depth of his feelings against the Vietnam War and the draft. He engaged in no violence, made no loud or unusual noise, uttered no sound before his arrest, and provoked no violent reaction or threat from anyone present.

\section{Procedural History}

Cohen was convicted in Los Angeles Municipal Court under California Penal Code \S~415, which barred ``maliciously and willfully disturb[ing] the peace or quiet of any neighborhood or person'' by ``offensive conduct,'' and was sentenced to 30 days. The California Court of Appeal affirmed, construing ``offensive conduct'' as ``behavior which has a tendency to provoke others to acts of violence or to in turn disturb the peace,'' and reasoning that someone might foreseeably strike Cohen or forcibly remove the jacket. After the California Supreme Court denied review, 4--3, the U.S.\ Supreme Court noted probable jurisdiction and reversed.

\section{Judicial Votes}

\textbf{5--4}, reversed (June 7, 1971).
\begin{itemize}
\item \textbf{Majority:} Justice Harlan, joined by Justices Douglas, Brennan, Stewart, and Marshall.
\item \textbf{Concurrences:} None.
\item \textbf{Dissent:} Justice Blackmun, joined by Chief Justice Burger and Justice Black, and as to its second paragraph by Justice White.
\end{itemize}

\section{Holding}

Absent a more particularized and compelling reason, a State may not make the simple public display of a single four-letter expletive a crime. Profanity used to convey political sentiment is protected speech, covering both cognitive and emotive expression, and passersby in a public courthouse corridor are not a captive audience: they can avoid offense by averting their eyes.

\section{Analysis of Opinions}

\textbf{Majority (Harlan): clearing the exceptions.} The conviction rested solely on the offensiveness of Cohen's words---pure speech, not separately identifiable conduct regulable under \textit{United States v.\ O'Brien}. Nor was \S~415 a courthouse-decorum rule, since it applied statewide and gave no notice of location-specific limits. The slogan was not obscene, for obscenity must be ``in some significant way, erotic.'' It was not fighting words under \textit{Chaplinsky}, being directed at the draft rather than ``to the person of the hearer.'' And there was no captive audience: outside the home, citizens must tolerate objectionable expression.

\textbf{Majority (Harlan): no power to cleanse vocabulary.} The violence rationale was ``undifferentiated fear or apprehension of disturbance'' (\textit{Tinker}). Harlan then set out three pillars: officials cannot draw principled lines among offensive words, so taste is left to the individual; language serves emotive as well as cognitive functions, and both are protected; and forbidding particular words risks suppressing ideas.

\textbf{Blackmun (dissenting).} Joined by Burger and Black, he called Cohen's ``absurd and immature antic'' ``mainly conduct, and little speech,'' within \textit{Chaplinsky}. Joined also by White, he would have remanded in light of \textit{In re Bushman}, which had narrowed \S~415 to a ``clear and present danger'' of violence.

\textbf{Later developments.} \textit{Iancu v.\ Brunetti} (2019) struck down the bar on ``immoral or scandalous'' trademarks, and \textit{Mahanoy v.\ B.L.} (2021) protected a cheerleader's profane off-campus Snapchat post.

\section{Examples: Future Applications}

\begin{enumerate}
\item \textbf{Same-Side: City Hall Sign.} A protester's profane sign condemning the mayor's housing policy is not fighting words or obscene, and passersby can look away; the arrest is unconstitutional.
\item \textbf{Same-Side: Bumper Stickers.} A ban on profane bumper stickers reaches emotive political speech by people who can avert their eyes, and gives officials the boundless power over words that \textit{Cohen} rejected.
\item \textbf{Opposite-Side: Face-to-Face Epithet.} Slurs hurled inches from a rival, daring him to fight, are \textit{Chaplinsky} fighting words, punishable under a narrowly construed statute.
\item \textbf{Opposite-Side: Broadcast Into the Home.} Afternoon radio profanity reaches homes and children who cannot simply look away; \textit{FCC v.\ Pacifica} (1978) upheld an indecency order.
\item \textbf{Fence-Sitter: Courtroom Apparel Ban.} A narrow ban on profane apparel inside a courtroom in session gives notice and governs a nonpublic forum, but invites arbitrary enforcement unless tied to disruption.
\end{enumerate}

\section{Critique}

\textbf{Civility.} Communitarian critics argue that Harlan's ``avert your eyes'' standard undervalues the shared interest in basic decorum, especially in solemn civic spaces like courthouses or around children, and treats every norm of restraint as a slippery slope to censorship.

\textbf{Line-drawing.} The Court later drew content-based lines around profanity in broadcasting (\textit{Pacifica}) and school assemblies (\textit{Bethel School District v.\ Fraser}), suggesting that \textit{Cohen}'s skepticism of principled distinctions is contextual rather than absolute.

\section{Key Quotations}

\begin{quote}
``For, while the particular four-letter word being litigated here is perhaps more distasteful than most others of its genre, it is nevertheless often true that one man's vulgarity is another's lyric. Indeed, we think it is largely because governmental officials cannot make principled distinctions in this area that the Constitution leaves matters of taste and style so largely to the individual.''
\end{quote}
\noindent --- Justice Harlan, for the Court, 403 U.S.\ at 25.

\begin{quote}
``Those in the Los Angeles courthouse could effectively avoid further bombardment of their sensibilities simply by averting their eyes.''
\end{quote}
\noindent --- Justice Harlan, for the Court, 403 U.S.\ at 21.

\begin{quote}
``In fact, words are often chosen as much for their emotive as their cognitive force. We cannot sanction the view that the Constitution, while solicitous of the cognitive content of individual speech, has little or no regard for that emotive function which, practically speaking, may often be the more important element of the overall message sought to be communicated.''
\end{quote}
\noindent --- Justice Harlan, for the Court, 403 U.S.\ at 26.

\end{document}
